% Research continuation of Differential Section 09,
% ``Higher spectral derivatives of the residue calculus.''
\section{Canonical spatial kernels for every spectral coefficient}
\label{log:section}

Let $P\in\mathbb C[u_0,\ldots,u_R]$ be fixed and write
$f_P(x)=P(\psi(x),\psi'(x),\ldots,\psi^{(R)}(x))$.
Its scalar projection is $\phi_P(Y)=P(Y,0,\ldots,0)$.
For a sufficiently large finite $M$, define the principal-part data
\[
 p_j(n)=[v^{-j}]f_P(-n+v),\qquad
 P_0(x)=\sum_{j=1}^M p_j(0)x^{-j},\qquad
 \Delta p_j(n)=p_j(n)-p_j(n-1).
\]
All $p_j(n)$ are explicit polynomials in ordinary harmonic numbers
and zeta values. Define the harmonic Dirichlet functions by
\begin{equation}\label{log:Dirichlet}
 D_j(s)=\sum_{n\ge1}\Delta p_j(n)n^{-s}\qquad(\Re s>0).
\end{equation}
The subscript $j$ denotes pole order, not differentiation order.
We are interested in every Laurent coefficient $[s^q]D_1(s)$,
$q\ge1$. The previously established unweighted finite-part bridge
uses only $q=1$ \cite{DPL,PTI}.

The normalized spatial function and its leading tail are
\begin{equation}\label{log:hK}
 h_P(x)=f_P(x+1)-f_P(x)+P_0(x),\qquad
 K_P(L)=\phi'_P(L)+p_1(0)=\sum_{k\ge0}K_kL^k.
\end{equation}
The $P_0$ term removes the entire singularity at $x=0$.
In particular $h_P(0)=f_P(1)-[x^0]f_P(x)$.

\begin{theorem}[Normalized Mellin identity]
\label{log:Mellin}
The integrals
\begin{equation}\label{log:E}
 E_P(s)=\int_0^1 x^{-s}h_P(x)\,dx
  +\int_1^\infty x^{-s}
       \left\{h_P(x)-\frac{K_P(\log x)}x\right\}\,dx
\end{equation}
define a holomorphic function on $-1<\Re s<1$.
On that strip, as an identity of meromorphic functions,
\begin{equation}\label{log:master}
 \boxed{\displaystyle
 D_1(s)=-\frac{\sin\pi s}{\pi}
       \left(E_P(s)+\sum_{k\ge0}\frac{K_k k!}{s^{k+1}}\right)
 -\sum_{j=2}^M\frac{(s)_{j-1}}{(j-1)!}D_j(j-1+s).}
\end{equation}
The rising factorials are finite polynomials.
Every shifted $D_j$ in the second sum is holomorphic at $s=0$,
with all its derivatives represented by absolutely convergent series.
\end{theorem}

\begin{proof}
The inherited normalized principal-part expansion \cite[Section 2]{DPL}
gives, on
$x>0$, the particularly simple shift identity
\begin{equation}\label{log:shiftres}
 h_P(x)=-\sum_{n\ge1}\sum_{j=1}^M
          \frac{\Delta p_j(n)}{(n+x)^j}.
\end{equation}
It can also be obtained directly by shifting a finite cutoff of the
principal-part expansion: the last boundary term tends to zero since
$p_j(n)=O((1+\log n)^d)$.
The local Laurent expansion proves analyticity at $x=0$.
The differentiable positive-axis polygamma asymptotic expansions give
\[
 f_P(x+1)-f_P(x)=\frac{\phi'_P(\log x)}x
                   +O((1+\log x)^d/x^2),
\]
and hence
$h_P(x)=K_P(\log x)/x+O((1+\log x)^d/x^2)$.
These estimates prove normal convergence of \eqref{log:E}, together
with every fixed $s$-derivative, on its stated strip.

For $0<\Re s<1$, integrate \eqref{log:shiftres} against $x^{-s}$.
The exact beta integral is
\[
 \int_0^\infty\frac{x^{-s}}{(n+x)^j}\,dx
 =n^{1-s-j}\frac{\Gamma(1-s)\Gamma(j-1+s)}{\Gamma(j)}
 =\frac{\pi}{\sin\pi s}\frac{(s)_{j-1}}{(j-1)!}
   n^{-(j-1+s)}.
\]
Fubini is justified by
$\Delta p_j(n)=O((1+\log n)^d/n)$; the slowest case, $j=1$,
is summable for $\Re s>0$ after taking the beta integral.
Also
\[
 \int_1^\infty x^{-s-1}K_P(\log x)\,dx
       =\sum_{k\ge0}K_k k!s^{-k-1}.
\]
Solving for $D_1(s)$ gives \eqref{log:master} on $0<\Re s<1$,
and its right side supplies continuation to the whole strip.
For $j\ge2$, $j-1+s$ has positive real part throughout
$\Re s>-1$, so its series is holomorphic there. This also proves
the assertions about ordinary derivatives.
\end{proof}

For $q\ge1$, define the universal polynomial
\begin{align}\label{log:kernel}
 k_q(L)&=\frac{(-1)^q}{q!}
          \frac{(L+i\pi)^q-(L-i\pi)^q}{2i\pi}\notag\\
 &=(-1)^q\sum_{m=0}^{\lfloor(q-1)/2\rfloor}
       \frac{(-1)^m\pi^{2m}L^{q-2m-1}}
            {(2m+1)!(q-2m-1)!}.
\end{align}
The first four are
\[
 k_1=-1,\qquad k_2=L,\qquad
 k_3=\frac{\pi^2}{6}-\frac{L^2}{2},\qquad
 k_4=\frac{L^3-\pi^2L}{6}.
\]
For any polynomial $w$, put
\[
 \mathcal L_P[w]=\int_0^1w(\log x)h_P(x)\,dx+
 \int_1^\infty w(\log x)
       \left\{h_P(x)-\frac{K_P(\log x)}x\right\}\,dx.
\]
Both integrals are ordinary absolutely convergent integrals.
Set
\begin{align}\label{log:constants}
 A_{P,q}&=-\sum_{\substack{k\ge0\\q+k\ \mathrm{even}}}
        \frac{(-1)^{(q+k)/2}\pi^{q+k}K_k k!}{(q+k+1)!},\notag\\
 T_{P,q}&=\sum_{j=2}^M\sum_{n\ge1}
       \frac{\Delta p_j(n)}{n^{j-1}}
       [s^q]\left\{\frac{(s)_{j-1}}{(j-1)!}e^{-s\log n}\right\}.
\end{align}
The inner coefficient is an explicit polynomial in $\log n$, so
$T_{P,q}$ is an absolutely convergent spectral sum. Equivalently it
is a finite rational combination of the derivatives
$D_j^{(r)}(j-1)$, $0\le r\le q-1$.

\begin{corollary}[Every Laurent coefficient has a normalized spatial kernel]
\label{log:coefficient}
For every fixed differential polynomial $P$ and every integer $q\ge1$,
\begin{equation}\label{log:coefficient-formula}
 \boxed{[s^q]D_1(s)=\mathcal L_P[k_q]+A_{P,q}-T_{P,q}.}
\end{equation}
All terms on the right are fixed by ordinary integrals, convergent
series, and finite endpoint/polynomial calculations.
\end{corollary}
\begin{proof}
Expand $\sin(\pi s)/\pi$ in \eqref{log:master} and differentiate
the normally convergent integrals coefficientwise. The coefficient
of $-\sin(\pi s)x^{-s}/\pi$ is exactly $k_q(\log x)$.
Multiplication with $K_k k!/s^{k+1}$ gives $A_{P,q}$ with the
stated sign and parity. The remaining finite sum gives $T_{P,q}$.
\end{proof}

\subsection{The same scalar projection controls every Stieltjes correction}

Define
\[
 R_P(L)=-\frac{\phi_P(L+i\pi)-\phi_P(L-i\pi)}{2i\pi},
 \qquad R'_P(L)=\sum_{r\ge0}a_rL^r.
\]
The inherited scalar-residue asymptotic \cite[Section 2]{DPL} is
\[
 \Delta p_1(n)=\frac{R'_P(\log n)}n
                 +O((1+\log n)^d/n^2).
\]
With $\zeta(1+s)=s^{-1}+\sum_{m\ge0}(-1)^m\gamma_m s^m/m!$,
the following all-order answer follows.

\begin{corollary}[Complete Stieltjes correction]
\label{log:Stieltjes}
For $q\ge1$,
\begin{equation}\label{log:stieltjes-formula}
 \boxed{\displaystyle
 [s^q]D_1(s)=\frac{(-1)^q}{q!}
 \left\{\sum_{r\ge0}a_r\gamma_{q+r}
 +\sum_{n\ge1}\left(\Delta p_1(n)-\frac{R'_P(\log n)}n\right)
                      (\log n)^q\right\}.}
\end{equation}
The second sum converges absolutely. In particular a mixed polynomial,
for which $\phi_P=0$, requires no Stieltjes correction at any order.
\end{corollary}
\begin{proof}
Subtract the leading polynomial from the Dirichlet series:
\[
 D_1(s)=\sum_r a_r(-1)^r\zeta^{(r)}(1+s)
  +\sum_{n\ge1}\left(\Delta p_1(n)-R'_P(\log n)/n\right)n^{-s}.
\]
The second series is holomorphic for $\Re s>-1$. The coefficient
of $s^q$ in $(-1)^r\zeta^{(r)}(1+s)$ is
$(-1)^q\gamma_{q+r}/q!$. This gives the formula.
\end{proof}

\paragraph{Recovery of the original finite-part normalization.}
For $q=1$, $T_{P,1}=\sum_{j\ge2,n\ge1}
\Delta p_j(n)/((j-1)n^{j-1})=J_P$.
If $C_P=(2\pi)^{-1}\int_{-\pi}^{\pi}\phi_P(i\theta)\,d\theta$,
then direct coefficient extraction gives
$A_{P,1}=\phi_P(0)-C_P$.
Thus the inherited identity
$\operatorname{FP}\int_0^1 f_P=C_P+J_P+[s]D_1(s)$ becomes
\begin{equation}\label{log:fpcheck}
 \operatorname{FP}\int_0^1f_P(x)\,dx=\phi_P(0)-E_P(0).
\end{equation}
This is a normalization check; no endpoint constant has been dropped.
For $q\ge2$, \eqref{log:coefficient-formula} supplies the spatial
normalization requested in Differential Section 09. The choice of the
splitting point $1$ and the full term $K_P(\log x)/x$ is part of the
definition. Altering it without transporting the elementary correction
would change the reported constants.

\subsection{Two explicit logarithmic integral evaluations}

For $P=u_0^2$, the pole differences are
$\Delta p_1(n)=-2/n$, $\Delta p_2(n)=0$, hence
$D_1(s)=-2\zeta(1+s)$ and $T_{P,q}=0$.
Moreover
\[
 h_P(x)=\frac{2(\psi(1+x)+\gamma)}x,\qquad
 K_P(L)=2L+2\gamma.
\]
The case $q=2$ gives the ordinary Stieltjes identity
\begin{equation}\label{log:gamma2}
 \boxed{\displaystyle
 \int_0^1\frac{\log x\,[\psi(1+x)+\gamma]}x\,dx
 +\int_1^\infty\frac{\log x\,[\psi(1+x)-\log x]}x\,dx
 =-\frac{\gamma_2}{2}-\frac{\pi^2\gamma}{6}.}
\end{equation}
The higher $q$ versions follow from the printed universal kernels.
These special-function ingredients have classical antecedents; the
claim here is the all-polynomial normalization theorem, not historical
priority for this isolated scalar integral.

For $P=u_1^2$, the principal parts are
\[
 p_4(n)=1,\quad p_2(n)=2(H_n^{(2)}+\zeta(2)),\quad
 p_1(n)=4(H_n^{(3)}-\zeta(3)),\quad p_3(n)=0.
\]
Thus $D_1(s)=4\zeta(3+s)$, $D_2(s)=2\zeta(2+s)$,
$T_{P,2}=2\zeta'(3)$, and $K_P=-4\zeta(3)$.
Write
\[
 h(x)=-\frac{2}{x^2}
        \bigl(\psi'(1+x)-\zeta(2)+2\zeta(3)x\bigr).
\]
The bracket has a double zero at zero. Equation
\eqref{log:coefficient-formula} gives
\begin{equation}\label{log:trigamma-square}
 \boxed{\displaystyle
 \int_0^1(\log x)h(x)\,dx
 +2\int_1^\infty\frac{\log x}{x^2}
          \bigl(\zeta(2)-\psi'(1+x)\bigr)\,dx
 =2\zeta''(3)+2\zeta'(3)+\frac{2\pi^2}{3}\zeta(3).}
\end{equation}
Both displayed integrals converge absolutely. This example explicitly
tests the origin subtraction, the higher-pole term $T_{P,2}$, and the
$\pi^2$ endpoint correction independently of the scalar Stieltjes case.
